\documentclass[11pt, a4paper]{article}
\usepackage[margin=2cm]{geometry}
\usepackage{amsmath, amsfonts, amssymb, xcolor, float, amsthm, mathtools, enumitem, cancel, soul, booktabs}
\usepackage[calcwidth]{titlesec}
\usepackage{ragged2e}

% \usepackage{graphicx}
% \graphicspath{ {./imgs/} }

\usepackage[hidelinks]{hyperref}
\hypersetup{
    colorlinks=true,
    citecolor=blue,
    filecolor=blue,
    linkcolor=blue,
    urlcolor=blue
}
\usepackage[nameinlink,capitalise,noabbrev]{cleveref}
\crefname{example}{Example}{Examples}
\crefname{equation}{Equation}{Equations}

\usepackage[T1]{fontenc}
\usepackage{minted}

\setminted{
    fontsize=\normalfont,
    style=emacs,
    linenos,
    breaklines,
    breakanywhere,
    autogobble,
    tabsize=4,
    frame=lines,
    framesep=1mm,
    numbersep=0.5em,      
}

% ── Python minted shortcuts ──────────────────────────────────
\newcommand{\py}[1]{\mintinline{python}|#1|}
\newcommand{\Rcode}[1]{\mintinline{R}|#1|}

% so verbatim tokenisation works correctly
\newminted[pyblock]{python}{}
\newminted[Rblock]{R}{}


% Optional: tweak line number look
\renewcommand\theFancyVerbLine{\sffamily\scriptsize\arabic{FancyVerbLine}}
\usepackage{fvextra}
\DefineVerbatimEnvironment{pyout}{Verbatim}{
    fontsize=\small,
    frame=lines,
    framesep=2mm
}

\newcommand{\R}{{\rm I\!R}}
\let\biconditional\leftrightarrow

%Formatting the sections
\renewcommand{\thesection}{\arabic{section}}
% \renewcommand{\thechapter}{\Roman{chapter}}
\setcounter{secnumdepth}{3}
\titleformat{\section}[hang]{\bfseries\Large}{\thesection}{1em}{}
\titleformat{\subsection}[hang]{\bfseries\filleft\large}{\itshape\thesubsection}{1em}{}[\vspace{-1.5ex}\rule{1.4\titlewidth}{1pt}]

\newcounter{example}[section] % resets each section (optional)
\renewcommand{\theexample}{\thesection.\arabic{example}} % e.g. 1.3.1 (optional)
\newcommand{\exampleparagraph}{%
  \refstepcounter{example}%
  \paragraph{Example \theexample}%
}

% Theoretical Statements Section (scrapped not useful)
% \titleclass{\statement}{straight}[\subsubsection]
% \newcounter{statement}
% \titleformat{\statement}[hang]{\bfseries}{\itshape\thestatement\quad--}{1em}{}[]
% \titlespacing{\statement}{1em}{1em}{1em}[1em]
% \renewcommand{\thestatement}{\arabic{statement}}

%Theoretical Statements Envs
%amsthm
\newtheoremstyle{break}%
    {}{}
    {}{\parindent}
    {\bfseries}{:}
    {\newline}{}
\theoremstyle{break}
\newtheorem{lem}{Lemma}
\newtheorem{cor}{Corollary}
\newtheorem{prop}{Proposition}
\newenvironment{lemma}[1]{\begin{lem}[#1]\leavevmode\vspace{-1.5\baselineskip}}{\end{lem}}
\newenvironment{corollary}[1]{\begin{cor}[#1]\leavevmode\vspace{-1.5\baselineskip}}{\end{cor}}
\newenvironment{proposition}[1]{\begin{prop}[#1]\leavevmode\vspace{-1.5\baselineskip}}{\end{prop}}

%For boxing environment
\usepackage{tcolorbox}
\tcbuselibrary{skins, breakable, most}
\newtcolorbox{quotebox}[1][]{%
  enhanced jigsaw,
  sharp corners,
  colback=white,
  borderline={1pt}{-2pt}{black},
  fontupper={\setlength{\parindent}{20pt}},
  #1
}
%Theorem box
\newtcolorbox{theorembox}[1]{
    breakable = true,
    title = #1,
    coltitle = black,
    colback = white,
    fonttitle = \bfseries,
    enhanced,
    sharp corners,
    boxrule = 0.5pt,
    colframe = black,
    attach boxed title to top left = {xshift = 15mm, yshift = -5.5mm},
    boxed title style={
        frame code={
            \path[draw = black, fill = black]
                ([yshift=-1mm,xshift=-1mm]frame.north west)
                arc[start angle=0,end angle=180,radius=1mm]
                ([yshift=-1mm,xshift=1mm]frame.north east)
                arc[start angle=180,end angle=0,radius=1mm];

            \path[draw = black, fill = white]
                ([xshift=-2mm]frame.north west) -- ([xshift=2mm]frame.north east)
                [rounded corners=1mm]-- ([xshift=1mm,yshift=-1mm]frame.north east)
                -- (frame.south east) -- (frame.south west)
                -- ([xshift=-1mm,yshift=-1mm]frame.north west)
                [sharp corners]-- cycle;
        },
        interior engine = empty,
        colback = white,
    },
    width = 0.95\linewidth,
    top = 10mm,
    left = 5mm,
    overlay = {

    }
}
% \newenvironment{theorembox}[1]
    % {\begin{center}\begin{theobox}{#1}}
    % {\end{theobox}\end{center}}

%Definition box
\newtcolorbox{definitionbox}{
    breakable = true,
    title = Definition,
    fonttitle = \bfseries,
    colback = white,
    width = 0.85\linewidth,
    coltitle = black,
    frame hidden,
    sharp corners,
    enhanced,
    attach boxed title to top left = {xshift = -0.5cm, yshift = -0.1mm},
    boxed title style = {
        colback = white,
        frame hidden,
    },
    overlay = {
        \node (A) at (title.south west) {};
        \node (B) at (title.south east) {};
        \node (C) at (frame.south west) {};
        \draw [semithick] (A) -- (B);
        \draw [semithick, double distance = 2pt] (C) -- (frame.north west);
    },
}
% \newenvironment{definitionbox}
    % {\begin{center}\begin{defbox}}
    % {\end{defbox}\end{center}}

\newtcolorbox{proofbox}{
    enhanced,
    breakable,
    colback=white,
    frame hidden,
    width=0.85\linewidth,
    left=2cm,                       % reserve space on the left
    overlay={
        % left vertical rule
        \draw[semithick] ([xshift=1.8cm]frame.north west) -- 
                        ([xshift=1.8cm]frame.south west);
        % "Proof" label at the top-left reserved area
        \node[anchor=north west,font=\itshape]
        at ([xshift=0.2cm,yshift=-0.2cm]frame.north west) {proof};
    },
}
% \newenvironment{proofbox}
    % {\begin{center}\begin{boxproof}}
    % {\end{boxproof}\end{center}}

\newtcolorbox{ProofIllustrationBox}{
    breakable = true,
    colback = white,
    coltitle = black,
    sharp corners,
    enhanced,
    %segmentation
    segmentation style = {
        solid,
        black,
    },
    %font
    fontlower = \sffamily\itshape\scriptsize,
    %format
    size = fbox,
    halign lower = center,
    halign upper = center,
}
\newtcolorbox{example}{
    breakable = true,
    title = Example,
    sidebyside,
    sidebyside align = top seam,
    enhanced,
    coltitle = black,
    colback = white,
    %Title
    attach boxed title to top left = {xshift = 1cm,},
    boxed title style = {
        colback = white,
        colframe = black,
        skin = enhancedfirst jigsaw,
        size = small,
        arc = 1mm,
        bottom = -1mm,
        boxrule = 0.8pt,
    },
    %Segmentation
    segmentation style = {
        thick,
        solid,
        shorten <= 10pt,
        shorten >= 10pt,
    },
    sharpish corners,
    colframe = black,
    boxrule = 0.8pt,
    width = 0.8\linewidth,
}
\newenvironment{examplebox}
    {\begin{center}\begin{example}}
    {\end{example}\end{center}}



%Tikz
\usepackage{tikz}
\usetikzlibrary{spy, intersections, calc, angles, quotes, fit, decorations.pathmorphing}
\tikzset{help lines/.style 2 args = {very thin, color = #1, step = #2}}
\tikzset{Hai lines/.style = {help lines = magenta}}

%For Pseudocode
\usepackage[ruled,titlenotnumbered]{algorithm2e}
\NoCaptionOfAlgo
\SetAlgoCaptionSeparator{}
\SetAlCapNameFnt{\large\scshape}
\DontPrintSemicolon
\SetAlgoLongEnd
\SetAlgoInsideSkip{smallskip}
\SetKwProg{Fn}{function}{:}{end function}
\SetKwComment{Comment}{$\triangleright$\ }{}
\usepackage{caption}
%Highlighting the pseudo-code in the algorithm environment using Tikz fit
\newcommand{\boxit}[3]{
    \tikz[remember picture,overlay] \node (A) {};\ignorespaces
    \tikz[remember picture,overlay]{
            \node[yshift=3pt,fill=#1,opacity=.25,
            fit={ ($(A)+(-0.1,0.3\baselineskip)$) ($(A)+(#2\textwidth,-{#3}\baselineskip - 0.25\baselineskip)$) } ] {};
        }\ignorespaces
}
\newcommand\mycommfont[1]{\footnotesize\ttfamily\itshape #1\unskip}
\SetCommentSty{mycommfont}
\setlength{\algomargin}{2em}
\setlength{\interspacetitleruled}{6pt}
\setlength\algoheightrule{1.5pt}


% \title{\texttt{Python} Notes}   
% \date{\today}
% \author{Tri Hai Huynh \\ \\
        % Undergraduate Student of The University \emph{of} Adelaide, \\ \\
        % Adelaide, South Australia}

\begin{document}
Since each hour has $60$ minutes therefore contains $12$ 5-minute intervals, starting from minute $0$ until minute $55$. Therefore, a 
set of interval is $H_k=\{0,5,10,\ldots,55\}$ where the interval value is $v_i$ for $i\in H_k$. Let define the hourly dispatch price as
\begin{equation}
    x_h=\frac{1}{12}\sum_{i\in H_k}v_i
\end{equation}
where $h$ is the index of the hour, not necessarily be in the 24 hour limit.

However, I have a population of hours in a whole year. I want to estimate that population mean, hence the population parameter is 
\begin{equation}
    \theta = \frac1N\sum_{h=1}^{N}x_h
\end{equation}

I will use stratified sampling, each stratum is block of hours in a day such as morning peak hours, working hours, evening 
peak hours, night hours, sleep hours to capture the characteristics of the data set. Therefore, there are 5 strata ($M=5$). I want to create a sample 
of arbitrary size $n$. So my population mean now is,
\begin{equation}
    \mu=\frac1N\sum_{d=1}^{M}\sum_{h=1}^{N_d}x_{dh}=\sum_{d=1}^{M}w_d\mu_d
\end{equation}

where $w_d=N_d/N$, $\mu_d=\frac{1}{N_d}\sum_{h=1}^{N_d}x_{dh}$. The latter summation is the weighted average of strata means $\mu_d$.

\section*{Estimation}
Let the mean of each stratum $d$ be $\hat{\mu}_d$. In stratified sampling, I choose $n_d$ units from a stratum. Let $s$ be the chosen units such that $s_1=s_2=\ldots=s_5$ which means that taking exactly same number of 
units in each stratum. Assume that I have some arbitrary stratum $d$, $s_d$ units of choice. Therefore, let the number of samples can be taken from the stratum be $K$ then $K=\binom{N_d}{s_d}$, where $N_d$ is the 
population of the stratum in my case the number of combination is extremely large thus rendering the simple random sampling (SRS) procedure in each stratum impractical. 

Instead of choosing all $K$ samples of the stratum, I can choose $\hat{K}$ such that $\hat{K}\ll K$. Now, the expected population parameter is 
\begin{equation}\label{eq:ExpectedPopulationParameter}
    \mathbb{E}(T)=\sum_{k=1}^{\hat{K}}\hat{\theta}_k\pi_k
\end{equation}
where $\hat{\theta}_k=\bar{x}_{d,k}$, the arithmetic mean of stratum $d$ of sample $k$. Let $S_{d,k}=\{i_1,i_2,\ldots,i_{s_d}\}$, where if $h\neq l$ then $i_h\neq i_l$ and $i_h\in\{1,2,\ldots,N_d\}$. This is a sample $k$
of stratum $d$. I have that 
\begin{equation}\label{eq:MeanOfStratumOfSample}
    \bar{x}_{d,k}=\frac{1}{s_d}\sum_{i\in S_{d,k}}x_i
\end{equation}
To calculate bunch of $\bar{x}_{d,k}$'s for $k=1,2,\ldots,\hat{K}$ and every stratum. In \Rcode{R}, I use $\hat{\mu}_h$ via $\bar{x}_{d,k}$ (a biased estimator).
\begin{Rblock}
    drawn <- lapply(strata, function(stratum) {
        stratum[sample(nrow(stratum), s_d, replace=FALSE), ]
    })
    # Mean of each sample k of each stratum.
    mu_d_hat <- sapply(drawn, function(d) mean(d$Mean)) #$
\end{Rblock}

Since this is SRS, then each sample has the same probability of being selected of $1/K$. From \cref{eq:ExpectedPopulationParameter} and \cref{eq:MeanOfStratumOfSample} hence
\begin{equation}
    \begin{aligned}
        \mathbb{E}(T)&=\sum_{k=1}^{\hat{K}}\bar{x}_{d,k}\frac{1}{K}\\
        \mathbb{E}(T)&=\sum_{k=1}^{\hat{K}}\sum_{i\in S_{d,k}}x_i\frac{1}{Ks_d} 
    \end{aligned}
\end{equation}

Furthermore, to estimate $\theta$, I need $\hat{\mu}$, which is 
\begin{equation}
    \hat{\mu}=\bar{x}_k=\sum_{d=1}^{M}w_d\bar{x}_{d,k}
\end{equation}
Hence in \Rcode{R} it is
\begin{Rblock}
    sum(weights*mu_h_hat)
\end{Rblock}

\subsection*{Expected Popluation Parameter}
Given that each unit has a chosen probability of $s_d/N_d$. In \Rcode{R}, it is not possible to generate the first $\hat{K}$ combinations so I choose randomly by looping $\hat{K}$ times. The expected 
number of times a unit appears is
\begin{equation}
    \mathbb{E}[U_i]=\hat{K}\frac{s_d}{N_d}
\end{equation}
I will consider two case since the appearance of units is random. Firstly,
\begin{equation}\label{eq:BiasDerive}
    \begin{aligned}
        \hat{K}\frac{s_d}{N_d}\sum_{i=1}^{N_d}x_i&\geq\sum_{k=1}^{\hat{K}}\sum_{i\in S_{d,k}}x_i\\
        \frac{\hat{K}}{K}\frac{1}{N_d}\sum_{i=1}^{N_d}x_i&\geq\frac{1}{Ks_d}\sum_{k=1}^{\hat{K}}\sum_{i\in S_{d,k}}x_i\\
        \frac{\hat{K}}{K}\mu_d&\geq\mathbb{E}(T)\\
        \frac{\hat{K}-K}{K}\mu_d&\geq\mathrm{bias}\\
    \end{aligned}
\end{equation}
For the other case it is 
\begin{equation}
    \frac{\hat{K}-K}{K}\mu_d<\mathrm{bias}
\end{equation}
Through the derivation, somehow the $s_d$ is eliminated but in \cref{eq:BiasDerive}. If I make $s_d$ somehow large then the bias just be $-\mu_d$. Illustrated with the graph below.

$\bar{x}_{d,k}$ remains to be a biased estimator of $\mu_d$. In addition to that, $w_d$'s are not all equal. Therefore, because of the previous two reasons, $\bar{x}_k$ is a biased estimator of $\mu$. However, by choosing
the correct amount of $s_d,\hat{K}$, the bias can be reduced. In \Rcode{R}, I use the \Rcode{replicate(K_hat, {...})} function to repeat the sampling procedure $\hat{K}$ times and calculate the mean of the whole sample
each for each sampling. The vector \Rcode{sample_means} contains the means of the sampling plans, that is, $\{\hat{\theta}_1,\hat{\theta}_2,\ldots,\hat{\theta}_{\hat{K}}\}$ of $\{S_1,S_2,\ldots,S_{\hat{K}}\}$ where 
each has a uniform probability of $\pi_k=1/K$ for $k=1,2,\ldots,\hat{K}$.
\begin{Rblock}
    sample_means <- replicate(K_hat, {
        drawn <- lapply(strata, function(stratum) {
            stratum[sample(nrow(stratum), s_d, replace=FALSE), ]
        })
        # Mean of each sample k of each stratum.
        mu_h_hat <- sapply(drawn, function(d) mean(d$Mean)) #$

        sum(weights * mu_h_hat)
    }) 
\end{Rblock}
\end{document}